% Main entry point. Overleaf: select main.tex and pdfLaTeX (TeX Live 2024 or newer).
% Compile with latexmk -pdf main.tex. Biber handles both reference sections.
% EDITING MAP: metadata.tex = identity; sections/ = manuscript content;
% references.bib = citations; ecology.sty = shared formatting.
% Keep this entry point mostly unchanged when writing your own manuscript.
\documentclass[11pt]{article}
\providecommand{\ReviewMode}{1} % 1 = review; 0 = reading. See reading.tex.
% EDIT THIS FILE FIRST. These values are shared by the manuscript and response.
% Replace every example name, affiliation, email and manuscript identifier.
% Use \& for a literal ampersand, \% for a percent sign and \_ for an underscore.
\newcommand{\ManuscriptTitle}{A clear question about ecological communities}
% ShortTitle appears in running headers; keep it brief.
\newcommand{\ShortTitle}{Ecological communities}
\newcommand{\ArticleType}{Research article}
% Journal and manuscript ID are used in the response letter, not a title-page box.
\newcommand{\TargetJournal}{Journal name}
\newcommand{\ManuscriptID}{Manuscript ID}
\newcommand{\RevisionDate}{4 October 2026}
% AuthorNames is plain text for PDF metadata and the supplement heading.
% Update AuthorBlock below as well; it is the formatted title-page author list.
\newcommand{\AuthorNames}{Jane Doe and Alex Smith}
\newcommand{\CorrespondingEmail}{jane.doe@example.org}
\newcommand{\Keywords}{community ecology; population dynamics; quantitative methods}

% Match each superscript to an affiliation. Add authors with the same pattern.
% The reading copy retains all identities; it does not anonymize a manuscript.
\newcommand{\AuthorBlock}{Jane Doe$^{1}$ and Alex Smith$^{2}$\\[.5em]
  \small $^{1}$Department of Ecology, Example University\\
  \small $^{2}$Department of Applied Mathematics, Example Institute\\
  \small Correspondence: \texttt{\CorrespondingEmail}}

\usepackage{ecology}
\addbibresource{references.bib}
% Optional counts remain available as macros but are not printed on the title page.
% Regenerate wordcounts.tex with count-words.py after editing your prose.
\InputIfFileExists{wordcounts.tex}{}{%
  \newcommand{\AbstractWordCount}{Not supplied}%
  \newcommand{\MainWordCount}{Not supplied}}
\title{\ManuscriptTitle}
\author{\AuthorBlock}
\date{}
\begin{document}
\maketitle
\thispagestyle{plain}
\begin{center}
  {\small\sffamily\ArticleType\quad\textbullet\quad\RevisionDate}
\end{center}

% A refsection gives this part an independent bibliography from the shared database.
% Keep citations and printbibliography inside the same refsection.
\begin{refsection}
% Start line numbering after the title; the reading entry point disables it.
\BeginReview
\section*{Abstract}
% Replace this paragraph with your abstract. Keep headings and keywords in main.tex.
% This file is counted separately by count-words.py.
This example demonstrates a writing workflow for quantitative ecology. We use a familiar population model to show equations, a mathematical statement, a figure and a table, with a separate supplement and bibliography. The text and illustrative results are teaching examples, not findings from a study. Replace them with the biological question, methods, evidence and conclusion of your manuscript.
\par
\smallskip\noindent\textbf{Keywords:} \Keywords
% Replace the illustrative prose and model with your own study.
% Keep stable, unique labels (sec:, eq:, fig:, tab:) and refer with \cref.
% Example: \label{fig:your-result} belongs AFTER its caption.
% TeXcount: count prose only; headers, captions and mathematics are excluded by the helper.
\section{Introduction}\label{sec:introduction}
A manuscript should make its biological question clear before introducing the machinery used to answer it. State what is known, what remains uncertain and why the proposed analysis can resolve that uncertainty. Classic population models offer a convenient example for demonstrating notation and references \citep{lotka1925,may1972}.

\linelabel{line:question-start}% Marks the first line of this paragraph in review mode.
The example question is how density-dependent growth approaches a carrying capacity. Every symbol, assumption and measurement should be defined before it carries an argument. The model below is illustrative; it is not presented as a fitted model or evidence about a real population.\linelabel{line:question-end}

\section{Methods}\label{sec:methods}
\subsection{Model and assumptions}
Let $N(t)>0$ denote abundance, $r>0$ the intrinsic growth rate and $K>0$ the carrying capacity. Logistic growth is written as
\begin{equation}\label{eq:logistic}
  \frac{\dd N}{\dd t}=rN\left(1-\frac{N}{K}\right).
\end{equation}
The same model can be written using scaled abundance $n=N/K$ and time $\tau=rt$:
\begin{align}
  \frac{\dd n}{\dd\tau} &= n(1-n),\label{eq:scaled}\\
  n(\tau) &= \frac{1}{1+(n_0^{-1}-1)e^{-\tau}}.\label{eq:solution}
\end{align}
Use \texttt{\textbackslash cref} for references such as \cref{eq:logistic,eq:scaled}; the object name is generated with the number. Use \texttt{\textbackslash citep} for parenthetical citations and \texttt{\textbackslash citet} for narrative citations. For example, \citet{may1972} connects ecological questions with mathematical analysis.

\begin{ecobox}{What the example assumes}
The environment and parameters are constant, abundance is continuous, and the population is closed. These assumptions determine the scope of the model. Explain their biological meaning before interpreting a result.
\end{ecobox}

\subsection{Mathematical statements}
\begin{proposition}\label{prop:limit}
For $n_0>0$, the solution in \cref{eq:solution} converges to $1$ as $\tau\to\infty$.
\end{proposition}
\begin{proof}
The denominator of \cref{eq:solution} tends to $1$ because $e^{-\tau}\to0$. It remains positive for $n_0>0$ and $\tau\geq0$, so the solution is well-defined.
\end{proof}
\Cref{prop:limit} separates a demonstrated mathematical conclusion from its biological interpretation. Additional derivations belong in \cref{sec:supp-derivation}.

\section{Results}\label{sec:results}
The illustrative trajectories in \cref{fig:trajectories} approach the same equilibrium from different initial abundances. They are sampled model trajectories, not observations, fitted estimates or uncertainty intervals. \Cref{tab:notation} keeps notation and units explicit.
% TeXcount ignores this float block. Put new figures and tables inside equivalent
% TC markers if they should be excluded from prose-only word counts.
% For your own plot, replace the input below with \includegraphics[width=\linewidth]{file}.
% Store the image in figures/ and keep the caption and label here.
%TC:ignore
\begin{figure}[htbp]
  \centering
  % Self-contained vector figure: no external image, font, shell escape or plotting runtime.
% Positions are samples of n(tau)=1/[1+(1/n0-1)exp(-tau)].
\setlength{\unitlength}{1mm}
\begin{picture}(130,66)(-12,-12)
\put(0,0){\vector(1,0){108}}\put(0,0){\vector(0,1){49}}
\multiput(0,45)(4,0){26}{\line(1,0){1.5}}
\put(-8,43){\small $1$}\put(-8,-2){\small $0$}
\put(101,-6){\small $6$}\put(46,-11){\small Scaled time $\tau$}
\put(-12,50){\small $n$}
\linethickness{.8pt}
\color{EcoAccent}
\qbezier(0,4.5)(9,7)(17,10.44)
\qbezier(17,10.44)(26,17)(34,20.28)
\qbezier(34,20.28)(43,28)(51,31.08)
\qbezier(51,31.08)(60,36.5)(68,38.61)
\qbezier(68,38.61)(77,42)(85,42.43)
\qbezier(85,42.43)(94,44)(102,44.02)
\color{EcoChange}
\qbezier[18](0,22.5)(9,29)(17,32.9)
\qbezier[18](17,32.9)(26,37)(34,39.64)
\qbezier[18](34,39.64)(43,42)(51,42.87)
\qbezier[18](51,42.87)(60,44)(68,44.19)
\qbezier[18](68,44.19)(77,44.6)(85,44.70)
\qbezier[18](85,44.70)(94,44.9)(102,44.89)
\color{black}
\put(4,51){\small Solid: $n_0=0.1$ \qquad Dashed: $n_0=0.5$}
\end{picture}

  \caption{Illustrative logistic trajectories. Scaled abundance is shown against scaled time for $n_0=0.1$ (solid line) and $n_0=0.5$ (dashed line). The dotted horizontal line marks the equilibrium $n=1$. Both trajectories follow \cref{eq:solution}; no empirical data are used.}
  \label{fig:trajectories}
\end{figure}
\begin{table}[htbp]
  \caption{Notation in the illustrative model. Units must be compatible with the abundance and time definitions.}
  \label{tab:notation}
  \centering
  \begin{tabularx}{\linewidth}{@{}lXl@{}}\toprule
    Symbol & Meaning & Units\\\midrule
    $N$ & Population abundance & Individuals\\
    $r$ & Intrinsic growth rate & Time$^{-1}$\\
    $K$ & Carrying capacity & Individuals\\
    $n$ & Scaled abundance, $N/K$ & Dimensionless\\
    $\tau$ & Scaled time, $rt$ & Dimensionless\\\bottomrule
  \end{tabularx}
\end{table}
%TC:endignore

% Flush pending Results floats before the Discussion; remove if your journal
% expects figures collected at the end instead.
\FloatBarrier
\section{Discussion}\label{sec:discussion}
Explain which conclusion follows from which evidence, distinguish mathematical results from empirical inference, and discuss assumptions that may fail in the system of interest. \revision{This sentence demonstrates revision highlighting: blue in the review copy, plain text in the reading copy.}

\section*{Acknowledgements}
Replace this example with acknowledgements and funding information.
\section*{Author contributions}
Specify who designed the study, developed the theory, collected or curated data, performed analyses and wrote the manuscript. Use the target journal's requested taxonomy.
\section*{Open research statement}
This demonstration uses no empirical data. Its figure follows the equation printed in the text. For your study, replace this paragraph with the actual data and code locations, access conditions, licenses and version identifiers.

\EndReview
\printbibliography[heading=bibintoc,title={References}]
\end{refsection}
% Resets counters and adds S prefixes; also starts the supplement-only contents.
% Do not reset these counters manually inside sections/supplement.tex.
\StartSupplement
\begin{refsection}
% Add supplementary sections here; S numbering is already configured by main.tex.
% Use unique labels (for example sec:supp-...) even when repeating a main-text topic.
% Cite normally: only references cited here enter the supplementary bibliography.
\section{Derivation of the solution}\label{sec:supp-derivation}
The supplement has its own section, equation, figure, table, theorem and page numbering. Its bibliography includes only references cited in this reference section. A reference can appear in both bibliographies when cited in both parts \citep{lotka1925}.

Separating variables in \cref{eq:scaled} yields
\begin{equation}\label{eq:supp-integral}
  \int\frac{\dd n}{n(1-n)}=\int\dd\tau.
\end{equation}
Integrating gives $\log[n/(1-n)]=\tau+C$ when $0<n<1$; the initial condition determines $C$. This produces \cref{eq:solution}. For $n_0>1$, an absolute value inside the logarithm yields the same final expression. The equilibrium case $n_0=1$ is handled directly.

\begin{lemma}\label{lemma:supp}
The scaled solution is constant when $n_0=1$.
\end{lemma}
\begin{proof}
Substitute $n_0=1$ into \cref{eq:solution}.
\end{proof}

\section{Reporting checks}\label{sec:supp-checks}
A reproducible analysis should identify data provenance, the model version, units, random seeds where relevant and the software environment. A concrete example of a simulation protocol for ecological models is described by \citet{grimm2006}.
%TC:ignore
\begin{table}[htbp]
\caption{Example supplement checklist. Replace each item with the information needed to reproduce your study.}\label{tab:supp}
\begin{tabularx}{\linewidth}{@{}lX@{}}\toprule
Item & What to record\\\midrule
Inputs & Source, version, access conditions and units\\
Model & Equations, assumptions and parameter values\\
Analysis & Code version, environment and numerical tolerances\\
Outputs & Figure scripts, intermediate products and uncertainty definitions\\\bottomrule
\end{tabularx}
\end{table}
\begin{figure}[htbp]
\centering% Self-contained vector figure: no external image, font, shell escape or plotting runtime.
% Positions are samples of n(tau)=1/[1+(1/n0-1)exp(-tau)].
\setlength{\unitlength}{1mm}
\begin{picture}(130,66)(-12,-12)
\put(0,0){\vector(1,0){108}}\put(0,0){\vector(0,1){49}}
\multiput(0,45)(4,0){26}{\line(1,0){1.5}}
\put(-8,43){\small $1$}\put(-8,-2){\small $0$}
\put(101,-6){\small $6$}\put(46,-11){\small Scaled time $\tau$}
\put(-12,50){\small $n$}
\linethickness{.8pt}
\color{EcoAccent}
\qbezier(0,4.5)(9,7)(17,10.44)
\qbezier(17,10.44)(26,17)(34,20.28)
\qbezier(34,20.28)(43,28)(51,31.08)
\qbezier(51,31.08)(60,36.5)(68,38.61)
\qbezier(68,38.61)(77,42)(85,42.43)
\qbezier(85,42.43)(94,44)(102,44.02)
\color{EcoChange}
\qbezier[18](0,22.5)(9,29)(17,32.9)
\qbezier[18](17,32.9)(26,37)(34,39.64)
\qbezier[18](34,39.64)(43,42)(51,42.87)
\qbezier[18](51,42.87)(60,44)(68,44.19)
\qbezier[18](68,44.19)(77,44.6)(85,44.70)
\qbezier[18](85,44.70)(94,44.9)(102,44.89)
\color{black}
\put(4,51){\small Solid: $n_0=0.1$ \qquad Dashed: $n_0=0.5$}
\end{picture}

\caption{A repeated demonstration to verify supplementary figure numbering. In a real manuscript, include additional evidence or diagnostics here rather than repeating a main-text figure.}\label{fig:supp}
\end{figure}
%TC:endignore

\EndReview
\printbibliography[heading=bibintoc,title={Supplementary references}]
\end{refsection}
\end{document}
